% !TEX root = ../../../main.tex
% 所属：第二部分：多变量分析学 / 第1章：Euclid空间
% 节：1.5 多元连续函数的性质
% 正文按原文件顺序保存；由 ../chapter1.tex 调用。
\section{多元连续函数的性质}
上一节已经定义了 $\R^n$ 中开集上映射的连续性. 下面定义一般的集合上的连续映射

\begin{definition}
	设点集 $K \subset \R^n$, $\vect f: K \to \R^m$ 为映射 , $\vect x_0 \in K$. 
	
	如果对于任意给定的 $\eps > 0$, 存在 $\delta > 0$, 使得当 $\vect x \in O(\vect x_0, \delta) \cap K$ 时, 成立
	\[
	|\vect f(\vect x) - \vect f(\vect x_0)| < \eps \quad (\text{即 } \vect f(\vect x) \in O(\vect f(\vect x_0), \eps)),
	\]
	则称 $\vect f$ 在点 $\vect x_0$ 连续.

	如果映射 $\vect f$ 在 $K$ 上每一点连续, 就称 $\vect f$ 在 $K$ 上连续, 或称映射 $\vect f$ 为 $K$ 上的连续映射.
\end{definition}

\begin{proposition}[连续映射将紧集映成紧集]
	\label{连续映射将紧集映成紧集}
	连续映射将紧集映成紧集.
\end{proposition}

\begin{proof}
	设 $K$ 是 $\R^n$ 中紧集, $\vect f: K \to \R^m$ 为连续映射. 要证明 $K$ 的像集
	\[
	\vect f(K) = \{ \vect y \in \R^m \mid \vect y = \vect f(\vect x),\, \vect x \in K \}
	\]
是紧集。由 Heine–Borel 定理及 Bolzano–Weierstrass 定理，Euclid 空间中的集合紧，当且仅当其中每个点列都有收敛子列，且子列极限仍属于该集合。因此，只需证明 $\vect f(K)$ 具有这一性质。

设 $\{\vect y_k\}\subset\vect f(K)$ 是任意点列。对每个 $k$，选取 $\vect x_k\in K$，使 $\vect f(\vect x_k)=\vect y_k$。由 $K$ 紧，存在子列 $\vect x_{k_i}\to\vect a\in K$。由连续性，
\[
\vect y_{k_i}=\vect f(\vect x_{k_i})\to\vect f(\vect a)\in\vect f(K).
\]
故 $\vect f(K)$ 是紧集。
\end{proof}

设 $f(\vect x)$ 是非空紧集 $K$ 上的连续函数, 那么 $f(K)$ 是 $\R$ 中的紧集, 因而有界闭, 且数集 $f(K)$ 取到最大数和最小数, 故:

\begin{theorem}[有界性定理]
	设 $K$ 是 $\R^n$ 中紧集, $f$ 是 $K$ 上的连续函数, 则 $f$ 在 $K$ 上有界.
\end{theorem}

\begin{proof}
	由定理 \ref{连续映射将紧集映成紧集}, $f(K)$ 紧, 故 $f(K)$ 在 $\R$ 中有界; $f(K)$ 即 $f$ 在 $K$ 上的值域, 故 $f$ 在 $K$ 上有界.
\end{proof}

\begin{theorem}[最值定理]
	设 $K$ 是 $\R^n$ 中非空紧集, $f$ 是 $K$ 上的连续函数, 则 $f$ 在 $K$ 上必能取到最大值和最小值, 即存在 $\xi_1, \xi_2 \in K$, 使得对一切 $\vect x \in K$ 成立
	\[
	f(\xi_1) \le f(\vect x) \le f(\xi_2).
	\]
\end{theorem}
\textbf{注：}这些定理的证明与 Part 1 第六章中对应定理的证明极其相似，这里略去（或者之后补充）

下面介绍一致连续性。
\begin{definition}
	设 $K$ 是 $\R^n$ 中点集, $\vect f: K \to \R^m$ 为映射. 
	
	如果对于任意给定的 $\eps > 0$, 存在 $\delta > 0$, 使得 $|\vect f(\vect x') - \vect f(\vect x'')| < \eps$ 对于 $K$ 中所有满足 $|\vect x' - \vect x''| < \delta$ 的 $\vect x', \vect x''$ 成立, 则称 $\vect f$ 在 $K$ 上一致连续.
\end{definition}
\begin{lemma}
	设 $K$ 是 $\R^n$ 中紧集, $\vect f: K \to \R^m$ 则 \( \vect f(\vect x) \) 在 \(K\) 上一致连续的充要条件是：
	\(\forall \{\vect x_n\}, \{\vect y_n\} \subset K, \lim\limits_{n \to \infty} |\vect x_n - \vect y_n| = 0\)，就有
	\[ \lim_{n \to \infty} |\vect f(\vect x_n) - \vect f(\vect y_n) |= 0 \]
	
	
\end{lemma}
\begin{proof}
	\textbf{必要性}：\( \vect f(\vect x) \) 在 \(K\)  上一致连续，\(\forall \varepsilon > 0, \exists \delta > 0, \) 使得 \(\forall \vect x_0, \vect y_0 \in K\)，
	且 \(|\vect x_0 - \vect y_0| < \delta\)，有 \(|\vect f(\vect x_0) - \vect f(\vect y_0)| < \varepsilon\).
	
	对上述\(\delta > 0\)，由 \(\lim\limits_{n \to \infty} |\vect x_n - \vect y_n|= 0\)，\(\exists N, \forall n > N, |\vect x_n - \vect y_n| < \delta\)，
	那么 \( |\vect  f(\vect x_n) -\vect f(\vect y_n)| < \varepsilon\).
	
	因此，\[\lim_{n \to \infty} |\vect  f(\vect x_n) -\vect f(\vect y_n)| = 0\]
\textbf{	充分性}：若不然，\(\vect f(\vect x)\) 在 \(K\) 上非一致连续：
	
	\(\exists \varepsilon_0 > 0, \forall \delta > 0, \exists \vect x_1, \vect x_2 \in K, \) 使得 \(|\vect x_1 - \vect x_2| < \delta\)，但 \(|\vect f(\vect x_1) - \vect f(\vect x_2)| \geqslant \varepsilon_0\).
	
	取 \(\delta_n = \dfrac{1}{n} (n = 1, 2, \dots)\)，于是有 \(\vect x_n, \vect y_n \in K\)，满足
	\[ |\vect x_n - \vect y_n| < \frac{1}{n}, \quad  |\vect  f(\vect x_n) -\vect f(\vect y_n)| \geqslant \varepsilon_0 \]
	
	显然，\(\lim\limits_{n \to \infty} |\vect x_n - \vect y_n| = 0\)，但 $|\vect f(\vect x_n)-\vect f(\vect y_n)|$ 不趋于 $0$，矛盾.
	
\end{proof}
上面引理常常用于证明一个映射不一致连续，当然也可以用于证明下面的定理，不过我们采取另一种方式。
\begin{theorem}[一致连续性定理]
	设 $K$ 是 $\R^n$ 中紧集, $\vect f: K \to \R^m$ 为连续映射, 则 $\vect f$ 在 $K$ 上一致连续.
\end{theorem}

\begin{proof}
	对于任意给定的 $\eps > 0$, 由 $\vect f$ 在 $K$ 上连续, 因此对任意 $\vect a \in K$, 存在 $\delta_{\vect a} > 0$, 使得当 $\vect x \in O(\vect a, \delta_{\vect a}) \cap K$ 时
	\[
	|\vect f(\vect x) - \vect f(\vect a)| < \frac{\eps}{2}.
	\]
	显然开集族 $\{ O(\vect a, \dfrac{\delta_{\vect a}}{2}): \vect a \in K \}$ 是 $K$ 的一个开覆盖. 由于 $K$ 是紧集, 因此存在其中有限个开集 $O(\vect a_1, \dfrac{\delta_{\vect a_1}}{2}), O(\vect a_2, \dfrac{\delta_{\vect a_2}}{2}), \cdots, O(\vect a_p, \dfrac{\delta_{\vect a_p}}{2})$ 覆盖了 $K$. 记 $\delta = \dfrac{1}{2} \min\{\delta_{\vect a_i}: 1 \le i \le p\}$. 那么对于 $K$ 中满足 $|\vect x' - \vect x''| < \delta$ 的任意 $\vect x'$ 和 $\vect x''$, 不妨设 $\vect x' \in O(\vect a_i, \dfrac{\delta_{\vect a_i}}{2})\ (1 \le i \le p)$, 则有
	\[
	|\vect x'' - \vect a_i| \le |\vect x'' - \vect x'| + |\vect x' - \vect a_i| < \frac{1}{2}\delta_{\vect a_i} + \frac{1}{2}\delta_{\vect a_i} = \delta_{\vect a_i},
	\]
	故 $|\vect f(\vect x'') - \vect f(\vect a_i)| < \dfrac{\eps}{2}$. 因此
	\[
	|\vect f(\vect x') - \vect f(\vect x'')| \le |\vect f(\vect x') - \vect f(\vect a_i)| + |\vect f(\vect x'') - \vect f(\vect a_i)| < \frac{\eps}{2} + \frac{\eps}{2} = \eps.
	\]
	由定义, $\vect f$ 在 $K$ 上一致连续.
\end{proof}
